\section{Валидность и границы обобщения}
\label{sec:heldout-limits}
Исследование проверяет один запрошенный псевдоним модели, один уровень рассуждения и одну версию CLI на выборке бенчмарка генерации функций. Повторы этапа разработки не дают репликации основной выборки на уровне моделей. Публичность задач оставляет нерешённым вопрос их присутствия в обучении; связанные задачи могут нарушать приближение независимых единиц. Другая модель, язык, распределение задач или версия сервиса способны изменить знак контраста. Даты и хеши среды идентифицируют исполненный клиент, но не неизменные веса провайдера.

Лимит времени и административное продолжение задают ещё одну границу. Неудачные генерации остаются пропусками, изменение правила исполнения раскрыто. Вывод по доступным парам относится к наблюдаемой оцениваемой части. Границы по всем назначениям показывают крайние варианты ненаблюдаемых исходов, но не исправляют смещение отбора. Превышение времени может отражать вычислительную сложность, а не случайный сетевой сбой, поэтому неотклонение нулевой разности качества не доказывает эквивалентность. Внешне обоснованные граница допустимой потери качества и совместный критерий качества и стоимости отсутствуют; решения о неухудшении или денежном превосходстве не принимаются.

Воздействия относятся к явным инструкциям. Названия ролей не удостоверяют наличие независимых агентов, а границы вызовов также раскрывают промежуточные ответы. Информация о задаче и операции этапов контролируются, но общий жёсткий лимит выходных токенов отсутствует. Ресурсные сравнения включают индуцированную длину рассуждений, повторную передачу контекста и накладные расходы клиента. Условная денежная оценка зависит от кеша; чувствительность без кеша убирает скидку, но не делает исполнение через подписку тождественным API-развёртыванию. Исследовательские вычисления и оценивание исключены; совокупная стоимость владения не оценивается.

Нативные тесты повышают проверяемость, сохраняя ограничения бенчмарка. Эталонный и неправильный контроли выявляют непригодную среду или нечувствительную отрицательную проверку, но их прохождение не доказывает соответствие каждого исходного утверждения естественно-языковой задаче. Отдельный аудит этапа разработки фиксирует расхождения инструкций и тестов как аннотации, сохраняя исходные промпты, тесты и оценки. Эти аннотации не служат исключениями после получения результатов или новым эталоном. Семь основных задач с неработающим эталоном также сохраняются в назначениях и учёте ресурсов.

Сравнение INoT относится к раскрытой промптовой адаптации с явным разрешением неоднозначности источника. Её отдельная серия сохраняет различия времени и кеша. Для Self-Collaboration выполнен лишь отдельный пилот авторской реализации на трёх задачах; его вариант без ролей и внешнее сравнение с достаточной мощностью не выполнены. Авторские реализации INoT, Agentless и SWE-agent также не оценивались. Последние системы добавляют поиск, инструменты и взаимодействие с репозиторием за пределами факторного вопроса с фиксированным контекстом. Проба SWE-bench на одном issue даёт настоящие патчи и тестовые исходы, но полная матрица десяти назначений на одной задаче не позволяет сообщать оценку репозиторного бенчмарка или рейтинг агентов. Для таких широких выводов нужны нативное сравнение на нескольких репозиториях и репликация на нескольких моделях при сопоставимом бюджете.
